\documentclass[10pt]{article}
% Shared LaTeX style for MIT 18.06 exam revision notes.
% Compile topic files with Tectonic, XeLaTeX, or LuaLaTeX.

\usepackage[a4paper,margin=0.72in]{geometry}
\usepackage{fontspec}
\usepackage{amsmath,amssymb,mathtools}
\usepackage{booktabs,array,tabularx}
\usepackage{enumitem}
\usepackage{titlesec}
\usepackage{fancyhdr}
\usepackage{microtype}
\usepackage{xcolor}

\setmainfont{Times New Roman}
\setsansfont{Arial}

\setlength{\parindent}{0pt}
\setlength{\parskip}{3.6pt}
\setlength{\abovedisplayskip}{6pt}
\setlength{\belowdisplayskip}{6pt}
\setlength{\abovedisplayshortskip}{4pt}
\setlength{\belowdisplayshortskip}{4pt}
\renewcommand{\arraystretch}{1.08}
\setlength{\tabcolsep}{4.5pt}

\titleformat{\section}
  {\normalfont\large\bfseries}
  {\thesection}
  {0.55em}
  {}
\titlespacing*{\section}{0pt}{10pt}{3pt}

\titleformat{\subsection}
  {\normalfont\normalsize\bfseries}
  {\thesubsection}
  {0.5em}
  {}
\titlespacing*{\subsection}{0pt}{7pt}{2pt}

\setlist[itemize]{leftmargin=1.35em,itemsep=1pt,topsep=2pt,parsep=0pt}
\setlist[enumerate]{leftmargin=1.55em,itemsep=1pt,topsep=2pt,parsep=0pt}

\pagestyle{fancy}
\fancyhf{}
\fancyfoot[C]{\scriptsize\color{black!45}MIT 18.06 Linear Algebra Exam Notes --- \thepage}
\renewcommand{\headrulewidth}{0pt}
\renewcommand{\footrulewidth}{0pt}

\newcommand{\notesubtitle}[1]{%
  {\centering\itshape #1\par}
  \vspace{2pt}
}

\newcommand{\source}[1]{%
  {\centering\small #1\par}
  \vspace{6pt}
}

\newcommand{\labelnote}[2]{%
  \par\smallskip
  \noindent\textbf{#1.}\quad #2
  \par\smallskip
}

\newcommand{\tightlabel}[1]{%
  \par\smallskip
  \noindent\textbf{#1.}\quad
}

\newcolumntype{Y}{>{\raggedright\arraybackslash}X}
\newcolumntype{L}[1]{>{\raggedright\arraybackslash}p{#1}}

\newenvironment{notetable}
  {\small\begin{tabularx}{\linewidth}}
  {\end{tabularx}}


\begin{document}

\begin{center}
{\LARGE Jordan Blocks and Defective Matrices\par}
\end{center}
\notesubtitle{Exam notes for repeated eigenvalues, missing eigenvectors, and Jordan form.}
\source{Based on handwritten MIT 18.06 revision notes.}

\section{Core Idea}

\labelnote{Core idea}{Repeated eigenvalues do not guarantee enough eigenvectors. When a matrix is missing eigenvectors, Jordan blocks replace diagonal form.}

Diagonalization is the best case:
\[
A=S\Lambda S^{-1}.
\]
This works only when \(A\) has a full basis of independent eigenvectors.

\section{Diagonalizable vs Defective}

\begin{center}
\small
\begin{tabularx}{\linewidth}{@{}L{0.30\linewidth}Y Y@{}}
\toprule
Case & Meaning & Exam signal \\
\midrule
Diagonalizable & Enough independent eigenvectors to form \(S\). & Each eigenspace supplies the needed number of vectors. \\
Defective & Not enough independent eigenvectors. & Some repeated eigenvalue has too small a nullspace. \\
Jordan form & Replacement for diagonal form. & Ones appear just above the diagonal inside a block. \\
\bottomrule
\end{tabularx}
\end{center}

\section{Algebraic and Geometric Multiplicity}

For an eigenvalue \(\lambda\):
\[
m_a(\lambda)=\text{multiplicity of }\lambda\text{ in }\det(A-\lambda I),
\]
\[
m_g(\lambda)=\dim N(A-\lambda I).
\]

The exam test is:
\[
A\text{ is diagonalizable}
\quad\Longleftrightarrow\quad
m_g(\lambda)=m_a(\lambda)\text{ for every }\lambda.
\]

If \(m_g(\lambda)<m_a(\lambda)\), the matrix is defective for that eigenvalue.

\section{Jordan Blocks}

A Jordan block for eigenvalue \(\lambda\) has \(\lambda\) on the diagonal and ones on the superdiagonal:
\[
J_{\lambda}=
\begin{bmatrix}
\lambda & 1 & 0\\
0 & \lambda & 1\\
0 & 0 & \lambda
\end{bmatrix}.
\]

The diagonal part gives the eigenvalue. The superdiagonal ones record the missing eigenvector directions.

\section{Jordan Form}

Instead of diagonalizing, a defective matrix may be written as
\[
A=MJM^{-1},
\]
where \(J\) is block diagonal with Jordan blocks.

A diagonal matrix is the special case where every Jordan block has size \(1\).

\section{Generalized Eigenvectors}

If \(x\) is an eigenvector, then
\[
(A-\lambda I)x=0.
\]

A generalized eigenvector \(w\) may satisfy
\[
(A-\lambda I)w=x.
\]

Together, \(w\) and \(x\) form a chain. This chain supplies the columns of \(M\) for the Jordan block.

\section{Exam Method}

\begin{enumerate}
  \item Find eigenvalues from \(\det(A-\lambda I)=0\).
  \item For each repeated eigenvalue, compute \(N(A-\lambda I)\).
  \item Compare \(m_g(\lambda)\) with \(m_a(\lambda)\).
  \item If all match, diagonalize.
  \item If one fails, use Jordan blocks for that eigenvalue.
\end{enumerate}

\section{Common Mistakes}

\begin{center}
\small
\begin{tabularx}{\linewidth}{@{}L{0.40\linewidth}Y@{}}
\toprule
Mistake & Correct exam habit \\
\midrule
Assuming repeated eigenvalues mean defective. & Check the dimension of the eigenspace. \\
Assuming distinct eigenvalues are required. & Repeated eigenvalues can still be diagonalizable. \\
Counting algebraic multiplicity as eigenvectors. & Count independent vectors in \(N(A-\lambda I)\). \\
Forgetting the superdiagonal ones in a Jordan block. & A nontrivial block has ones above the diagonal. \\
Treating \(A=MJM^{-1}\) as diagonalization. & Jordan form is weaker than diagonal form. \\
\bottomrule
\end{tabularx}
\end{center}

\section{Final Exam Checklist}

\tightlabel{Checklist}
\begin{enumerate}
  \item Compute the characteristic polynomial.
  \item List algebraic multiplicities.
  \item Compute eigenspace dimensions.
  \item Compare \(m_a(\lambda)\) and \(m_g(\lambda)\).
  \item Use diagonal form if enough eigenvectors exist.
  \item Use Jordan blocks when eigenvectors are missing.
  \item Check that the total block sizes add to \(n\).
\end{enumerate}

\end{document}
